% SPDX-License-Identifier: Apache-2.0
\section{A Program That Moves Beats Through a Wide Port}
\label{sec:x-stream}

\code{HostLink} and \code{LiteMaster} reach a 32-bit AXI-Lite bus, one
word per frame. FireSim's streams need more. Its host driver moves
them with \code{cpu\_managed\_axi4\_read} and
\code{cpu\_managed\_axi4\_write}, on an AXI4 port 128 bits wide with
64-bit addresses. This example reaches such a port, one 16-byte beat
per frame. \code{StreamLink} takes the ask off the wire,
\code{BeatMaster} performs it as an AXI4 burst of one beat, and the
answer goes back in a frame.

The frames have their own EtherType, \code{0x88b6}, next to the
\code{0x88b5} of \code{HostLink}'s frames. So \code{EthShare} can give
each unit its own frames by the type alone, and both can share the
board's one Ethernet port.

Behind the port is a queue, as a FireSim stream is: a read takes the
next beat out, and a write puts one in. So a transfer performed twice
loses a beat or adds one. As \code{LiteMaster} does, \code{BeatMaster}
keeps the tag and the answer of the last transfer, the beat it read
included, and answers an ask that repeats the tag from them.

A program starts with a tag the board may have kept from an earlier
run, and no stream transfer is safe to skip or to repeat. So its
first ask is a no-op, which \code{BeatMaster} answers without the
port and whose tag it keeps, whatever it kept before.

The run starts with that no-op. Then it writes two beats and reads
them back. It asks the second read
again under the same tag, as the program does when an answer is lost,
and the count of transfers the port saw does not move. Then it reads
the empty queue, and the error comes back in the answer.

The two units and \code{StreamBridge}, which joins them, are lowered,
and the build simulates the three netlists against this run under
\code{nvc} and Verilator.

\begin{figure}[htbp]
\centering
\input{streamlink_timing.tex}
\caption{The seven asks of the run. Each arrives as bytes on \code{rx}
and becomes one \code{ask}; \code{busy} is high while its transfer is
on the port, a read's address on \code{ar}, and the \code{back} it
brings leaves as a frame on \code{tx}. The first, the no-op, and the
sixth, the repeated read, are answered without \code{busy} or
\code{ar}.}
\label{fig:x-stream}
\end{figure}

\lstinputlisting[caption={\code{ex\_stream.rs}. Run with \code{bazel run //lib/examples:ex\_stream}.},label={lst:x-stream}]{lib/examples/ex_stream.rs}

\lstinputlisting[style=ascii,caption={What \code{ex\_stream} printed when the build ran it: each ask, its tag, the beat in its answer, and how many transfers the port had seen.},label={lst:x-stream-out}]{out_streamlink.txt}
